\documentclass[a4paper]{article}
\usepackage[polish]{babel}
\usepackage[cp1250]{inputenc}
\usepackage[T1]{fontenc}
\usepackage[dvips]{graphicx}

\pagestyle{headings}
\textwidth      15.5cm
\oddsidemargin    .1cm
\evensidemargin   .1cm

\begin{document}

\thispagestyle{empty}

\begin{minipage}{5cm}
\includegraphics[scale=0.7]{salomon.eps}
\end{minipage}
\begin{minipage}{10cm}
\begin{center}
{\Large Salomon\\
System przetwarzania wiedzy\\}
\end{center}
\end{minipage}

\vspace*{0.5cm}

\hrulefill

\vspace*{1cm}


\begin{flushright}
{\Large \emph{Iteration Description --- C1, ver 1.0} }
\end{flushright}



\begin{flushleft}
{ \Large Historia wersji }\[\]
\begin{tabular}{|p{2.5cm}||p{1.5cm}|p{6cm}|p{3.5cm}|}
  \hline
  Data & Wersja & Opis & Autor \\
  \hline \hline
  29.09.2005 & 1.0 & utworzenie & \emph{Leszek Grzanka}\\
  \hline
  \hline
\end{tabular}
\end{flushleft}

\newpage


\section{Zadania dla iteracji } 
Stabilizacja Salomona. Dostarczenie wersji realizuj�cej wymagania stawiane przez grup� wsp�tworz�c� Salomona. Implementacja warstwy \textit{Solution}.

\section{Produkt iteracji} 
Stabilna wersja Salomona, zawieraj�ca w pe�ni sprawn� warstw� \textit{Solution}, oraz mo�liwo�� pod��czania si� do zewn�trznych baz danych.

\section{Plan test�w} 


\section{Efekt wykonania iteracji}
Efektem by�o:
\begin{itemize}
\item zaimplementowanie odpowiednich klas realizuj�cych funkcjonalno�� warstwy \textit{Solution}
\item rozbudowa bazy danych o tabele przechwuj�c� informacj� o obiektach \textit{Solution}
\item modyfikacja interfejsu u�ytkownika o elementy umo�liwiaj�ce zarz�dzanie warstw� \textit{Solution}
\end{itemize}

\end{document}

